\def\topic{Polynomial trees with a class of labels}
\ifdefined\MyHeaderIncluded
\else
\def\qwertyflag{true}

\ifdefined\MyHeaderIncluded

\else
  \newcommand{\foobar}{FooBar}
  \def\MyHeaderIncluded{true}
% if is till the end of the file

\documentclass[10pt,a4paper]{article}
\usepackage[utf8]{inputenc}
%\usepackage[german]{babel}
%\usepackage[cyrillic]{babel}
\usepackage[T1]{fontenc}
%\usepackage[T2A]{fontenc}
%\usepackage[english, russian]{babel}
\usepackage{amsmath}
\usepackage{amsfonts}
\usepackage{amssymb}
\usepackage{amsthm}
\usepackage{mathtools} %coloneqq

\usepackage{bussproofs} %

%\renewcommand{\thesubsection}{\arabic{subsection}}
%\newtheorem{thm}{Theorem}[section]
%\renewcommand{\thethm}{\arabic{subsection}.\arabic{thm}}

\usepackage{thmtools}
\usepackage{hyperref}
\hypersetup{
    colorlinks=true,
    linkcolor=blue,
    filecolor=magenta,      
    urlcolor=cyan,
    pdftitle={Overleaf Example},
    %pdfpagemode=FullScreen,
    }
%\usepackage{stix-type1}%\turnediota
\theoremstyle{definition}

\newtheorem{constructor}{Constructor}[section]
\newtheorem{iaxiom}{Axiom}
\newtheorem{axiom}{Axiom}
\newtheorem*{axiom*}{Axiom}
\newtheorem{definition}{Definition}[section]
\newtheorem*{definition*}{Definition}
\newtheorem{remark}{Remark}[section]
\newtheorem*{remark*}{Remark}
\newtheorem{recall}{Recall}[section]
\newtheorem*{recall*}{Recall}

\newtheorem{lemma}{Lemma}[section]
\newtheorem{theorem}{Theorem}[section]
\newtheorem*{theorem*}{Theorem}
\newtheorem{note}{Note}[section]
\usepackage{etoolbox}

%\newcommand*{\@Dashv}[2]{%
%  \reflectbox{$\m@th#1#2$}%
%}
\newcommand*{\separation}[3]{% 1)binder, 2)set, 3)condition
  \{#1\in #2: #3\}
}

% https://www.freecodecamp.org/news/simple-regex-tricks-for-beginners-3acb3fa257cb/
%  \\mathcal\{T\} |---> \topology
%  \\mathcal\{([^\}]*)\}   |--->  $1
\newcommand{\AND}{\ \ \&\ \ }
\newcommand{\any}{\mathrm{\ any}}
\newcommand{\borel}{\mathcal{B}}
\newcommand{\can}{\mathrm{can}}
\newcommand{\canonize}{\mathrm{canonize}}
\newcommand{\closure}{\mbox{cl}}
\newcommand{\complex}{\mathbb{C}}
\newcommand{\converges}{\longrightarrow}
\newcommand{\countable}{\mbox{\ countable}}
\newcommand{\ldefined}{\mbox{\ defined}}
\newcommand{\dom}{\mbox{dom}}
\newcommand{\Dom}{\mbox{Dom}}
\newcommand{\Done}{\noindent \textbf{Done.}\\}
\newcommand{\lexists}{\mbox{\ exists}}
\newcommand{\lforall}{\mbox{\ forall}}
\newcommand{\lunique}{\mbox{\ unique}}
\newcommand{\family}{\mathrm{\ family}}
\newcommand{\Fm}{\mathrm{Fm}}
\newcommand{\FuncSymb}{\mathrm{FuncSymb}}
\newcommand{\function}{\mathrm{\ function}}
\newcommand{\functional}{\mbox{\ functional}}
\newcommand{\IFF}{\ \Longleftrightarrow\ }
\renewcommand{\iff}{\ \longleftrightarrow\ }
\newcommand{\IMPLIES}{\ \Longrightarrow\ }
\renewcommand{\implies}{\ \longrightarrow\ }
\newcommand{\Inf}{\mathrm{Inf}}
\newcommand{\injective}{\mbox{\ injective}}
\newcommand{\injection}{\mbox{\ injection}}
\newcommand{\interior}{\mbox{int}}
\newcommand{\intersection}{\bigcap}
\newcommand{\isempty}{\mathrm{\ empty}}
\newcommand{\LB}{\mathrm{LB}}
\newcommand{\nghbset}{\mathcal{N}}
\newcommand{\nonempty}{\mathrm{\ nonempty}}
\newcommand{\Ord}{\mathrm{Ord}}
\newcommand{\power}{\mathcal{P}}
\newcommand{\pr}{\mathrm{pr}}
\newcommand{\fst}{\mathrm{fst}}
\newcommand{\snd}{\mathrm{snd}}
\newcommand{\PredSymb}{\mathrm{PredSymb}}
\newcommand{\Proof}{\noindent \textbf{Proof.}\\}
\newcommand{\ran}{\mathrm{ran}}
\newcommand{\Ran}{\mathrm{Ran}}
\newcommand{\real}{\mathbb{R}}
\newcommand{\relation}{\mathrm{\ relation}}
\newcommand{\set}{\mathrm{\ set}}
\newcommand{\singleton}{\mathrm{\ singleton}}
\newcommand{\semantics}{\mathrm{semantics}}
\newcommand{\semialgebra}{\mathcal{S}}
\newcommand{\Sup}{\mathrm{Sup}}
\newcommand{\surjective}{\mbox{\ surjective\ on\ }}
\newcommand{\Term}{\mathrm{Term}}
\newcommand{\Type}{\mathrm{Type}}
\newcommand{\topology}{\mathcal{T}}
\newcommand{\total}{\mbox{\ total\ on\ }}
\newcommand{\transitive}{\mathrm{\ transitive}}
\newcommand{\union}{\bigcup}
\newcommand{\lTrue}{\mathrm{True}}
\newcommand{\lFalse}{\mathrm{False}}
\newcommand{\lUncertain}{\mathrm{Uncertain}}
\newcommand{\lif}{\mathrm{if}}
\newcommand{\lthen}{\mathrm{then}}
\newcommand{\lelse}{\mathrm{else}}
\newcommand{\lexception}{\mathrm{exception}}
\newcommand{\lundef}{\mathrm{undef}}
\newcommand{\univocal}{\mathrm{\ univocal}}
\newcommand{\univocalat}{\mathrm{\ univocal\ at\ }}
\newcommand{\UB}{\mathrm{UB}}
\newcommand{\Un}{\mathcal{U}}
\newcommand{\Var}{\mathrm{Var}}

\newcommand{\notvarphi}{\phi}
\renewcommand{\phi}{\varphi}

\newcommand{\comment}[1]{\quad\mbox{(#1)}}
% local commands
%\preto\section{\undef\blipo} %etoolbox
\preto\subsection{\undef\blipo} %etoolbox
% https://tex.stackexchange.com/questions/105684/limit-macro-newcommand-scope-to-section
\usepackage{graphicx}
\usepackage{authblk}

\setlength\parindent{0pt} % whole doc of \noindent
\usepackage{soul}
%\usepackage{mathgl}

\fi % guard
\ifdefined\topic
\else
\def\topic{No Titel}
\fi
\title{\topic}
\author[1]{Georgy Dunaev}
\begin{document}
\maketitle
\fi
\section{Polynomial trees with a class of labels}
\label{sec:class_polynomial_trees}

The finite-address construction extends to a proper class of labels when
every label is a set and every branch fibre is a set. The carriers may
then be proper classes, although each individual tree remains a set.
This gives an intermediate model for the larger class-valued proposal
in the preceding section.

\subsection{Local assumptions and class operations}
Work in classical NBG with the ordinary set axioms, including Pairing,
Union, Separation, Infinity, powerset, and Replacement for set domains
with fixed class parameters; MK also suffices. Only elementary
class comprehension is used: its quantified variables range over sets,
and previously specified classes occur as parameters. No Choice or
Global Choice is used. Auxiliary pairs and function graphs use ordinary
Kuratowski coding. The founded-recursion arguments below use
Theorem~\ref{thm:setlike_class_recursion}; foundedness is proved for the
particular relations, so global regularity is not needed for those steps.

Fix a class $A$ and a class relation $R\subseteq A\times\V$ such that
every fibre
\[
B(a)=\{b:(a,b)\in R\}\qquad(a\in A)
\]
is a set. Here $B(a)$ denotes this specified fibre, independently of the
historical family encodings. Its uniqueness follows from extensionality;
no choice of a representative is involved. Since labels belong to $A$,
they are sets even when $A$ is proper.

For any class $X$, define
\[
F_B(X)=\{(a,s):a\in A,\ s\text{ is a set function }B(a)\to X\}.
\]
This is a class of set pair codes. For a class function $u:X\to Y$,
Replacement on the set $B(a)$ makes $u\circ s$ a set graph, and
\[
F_B(u)(a,s)=(a,u\circ s)
\]
defines a class function $F_B(X)\to F_B(Y)$. The equations
$F_B(1_X)=1_{F_B(X)}$ and
$F_B(v\circ u)=F_B(v)\circ F_B(u)$ follow by evaluation at each branch.

The universal properties below mean statements for every specified
class $X$ and class structure map, with existence and uniqueness of a
class map satisfying the stated equation. They do not form a class
whose members are all classes or all class-function graphs. This
interpretation avoids needing an internal category of all classes.

\subsection{Finite addresses and the size of an individual tree}
Put
\[
C=\{(a,b):a\in A,\ b\in B(a)\}.
\]
A finite word over $C$ is a set function $p:n\to C$ for some
$n\in\omega$. Write $\operatorname{Word}(C)$ for the class of these
words, $\epsilon$ for the empty word, $pq$ for concatenation, and
$pc$ for appending one letter $c$. Every such word is a set, although
$C$ and $\operatorname{Word}(C)$ may be proper classes.

\begin{definition}[Class-labelled finite-address tree]
A tree is a functional class graph $t:D_t\to A$, with
$D_t\subseteq\operatorname{Word}(C)$, satisfying:
\begin{enumerate}
\item $\epsilon\in D_t$;
\item $D_t$ is prefix closed;
\item for $p\in D_t$ and $c\in C$,
\[
pc\in D_t\iff c=(t(p),b)\text{ for some }b\in B(t(p)).
\]
\end{enumerate}
The functional graph has set words and set labels as its coordinates.
\end{definition}

\begin{lemma}[Every such individual tree is a set]
\label{lem:class_tree_is_set}
A class-labelled finite-address tree has a set domain and a set graph.
\end{lemma}
\Proof
For each $n\in\omega$, let $D_n$ consist of its addresses of length
$n$. We prove that $D_n$ is a set by induction. The root condition
gives $D_0=\{\epsilon\}$. If $D_n$ is a set, then for each
$p\in D_n$ the exact-child condition gives the set of children
\[
\{p(t(p),b):b\in B(t(p))\}.
\]
It is a set by Replacement on the set fibre. Replacement over $D_n$
collects these child sets, and their union is $D_{n+1}$: every
nonempty address has its unique preceding prefix, which belongs to
$D_t$ by prefix closure. This proves the induction step.

The level $D_n$ is uniquely specified from $n$ and the fixed class
parameter $t$. Replacement over $\omega$ therefore collects the levels,
and their union is $D_t$, since every word has finite length.
Replacement on this set domain now collects the pairs $(p,t(p))$,
giving exactly the graph $t$ as a set. There is no selection of a
branch or of a label: both are supplied by the given graph.
\Done

Define $\mathcal M_B$ to be the class of set graphs satisfying these
tree conditions. This is an elementary comprehension instance: words,
their domains, and their lengths are sets, and $A,R$ are fixed class
parameters. Lemma~\ref{lem:class_tree_is_set} shows that it includes
every functional class tree with this finite-address interface.

\begin{lemma}[Class-labelled grafting and decomposition]
\label{lem:class_tree_grafting}
There is a class bijection
\[
\delta:\mathcal M_B\longrightarrow F_B(\mathcal M_B),\qquad
\delta(t)=(a,b\mapsto t_b),\quad a=t(\epsilon),
\]
where $t_b$ is obtained by removing the initial letter $(a,b)$ from
the addresses beginning with that letter. Its inverse is grafting.
\end{lemma}
\Proof
For $b\in B(a)$, Separation and Replacement on the set graph $t$
give the subtree with domain and values
\[
D_{t_b}=\{q:(a,b)q\in D_t\},\qquad t_b(q)=t((a,b)q).
\]
The tree conditions hold after removing the prefix. Replacement on
the set $B(a)$ then gives the set function $b\mapsto t_b$.

Conversely, given $(a,s)\in F_B(\mathcal M_B)$, use the domain
\[
D=\{\epsilon\}\cup
\{(a,b)q:b\in B(a),\ q\in D_{s(b)}\}.
\]
Replacement collects the prefixed subtree domains over the set $B(a)$;
union makes $D$ a set. Give the root label $a$ and give $(a,b)q$
the label $s(b)(q)$. Unique decomposition of a nonempty word into its
first letter and suffix makes this a well-defined set graph. Prefix
closure and the exact-child condition hold at the root by $B(a)$ and
below it by the corresponding conditions on $s(b)$.
Removing and restoring the first letter recovers every address and
label. Thus these operations are inverse. Their uniquely specified
set outputs define class maps by elementary comprehension.
\Done

\subsection{Corecursion for class state spaces}
\begin{theorem}[Finality for class coalgebras with set arities]
\label{thm:class_M_final}
For every class $X$ and class function $\gamma:X\to F_B(X)$, there
is exactly one class function $h:X\to\mathcal M_B$ satisfying
\[
\delta\circ h=F_B(h)\circ\gamma.
\]
\end{theorem}
\Proof
Write $\gamma(y)=(a_y,s_y)$. Fix $x\in X$. Construct the reached
address/state pairs by natural-number recursion:
\[
L_0=\{(\epsilon,x)\},\qquad
L_{n+1}=\{(p(a_y,b),s_y(b)):(p,y)\in L_n,\ b\in B(a_y)\}.
\]
Each successor is a set: for each record of the set $L_n$, Replacement
on $B(a_y)$ forms its next records; Replacement over $L_n$ and union
collect them. The operation is uniquely specified with fixed class
parameter $\gamma$, so ordinary set-valued recursion on $\omega$
applies. It can be made total on other set inputs by returning the
empty set there. Inductively every reached state belongs to $X$, so
that default is never used. Replacement over $\omega$ and union give
the set $L=\bigcup_nL_n$.

At each length an address reaches at most one state: remove its last
letter and use the uniqueness of the preceding state and of the
transition $s_y(b)$. Different lengths have disjoint address domains.
Thus $L$ is functional. Replacing each reached state $y$ by its label
$a_y$ gives a set graph $t_x$. Its domain contains the root, is prefix
closed, and has exactly the children prescribed by the reached state's
arity. Consequently $t_x\in\mathcal M_B$.

Its root has label $a_x$. After the first letter $(a_x,b)$ the same
transition construction starts from $s_x(b)$, so its subtree is
$t_{s_x(b)}$. This proves the equation for $h(x)=t_x$.
The construction uniquely specifies a set graph for each $x$ using
set quantifiers and the fixed parameters. Elementary class comprehension
therefore gives the class function $h$, even when $X$ is proper.

If $k:X\to\mathcal M_B$ also satisfies the equation, its root at $x$
has label $a_x$ and its $b$-subtree is $k(s_x(b))$. Induction on word
length, simultaneously for all $x\in X$, gives the same admissible
addresses and labels as the transition construction. Every address is
finite, so the graphs agree at each $x$, and $k=h$.
\Done

\begin{corollary}[Class bisimulation]
If a class relation $Q\subseteq\mathcal M_B\times\mathcal M_B$
relates only trees with equal root labels and relates their corresponding
subtrees at every branch, then $t\mathrel Q u$ implies $t=u$.
\end{corollary}
\Proof
For each related pair with common root label $a$, Replacement on $B(a)$
forms the set function $b\mapsto(t_b,u_b)$. It defines a class
coalgebra $Q\to F_B(Q)$. Both projections to $\mathcal M_B$ satisfy
the coalgebra-morphism equation. Theorem~\ref{thm:class_M_final} makes
them equal, proving the assertion.
\Done

\subsection{The founded class and its recursion principle}
Let $\mathcal W_B\subseteq\mathcal M_B$ consist of trees whose
immediate-child relation on $D_t$ is well founded, with a child preceding
its parent. This means every nonempty subset of $D_t$ has a minimal
member. The subsets quantified over are sets, so elementary comprehension
defines $\mathcal W_B$. This is the minimal-member condition, not an
assertion that no infinite branch has been chosen.

\begin{lemma}[Founded grafting and the subtree relation]
\label{lem:class_W_founded}
Grafting restricts to a class bijection
$F_B(\mathcal W_B)\to\mathcal W_B$. The immediate-root-subtree
relation on $\mathcal W_B$ is founded and set-like.
\end{lemma}
\Proof
For $t=\delta^{-1}(a,s)$, foundedness of $t$ implies foundedness of
each $s(b)$ by prefixing a nonempty subset of subtree addresses with
$(a,b)$ and taking a minimal member. Conversely, suppose every $s(b)$
is founded and take a nonempty subset $S$ of $D_t$. If $S$ contains
only the root, the root is minimal. Otherwise take a nonroot address
in $S$ with first letter $(a,b)$. The suffixes of the addresses of $S$
with that first letter form a nonempty subset of $D_{s(b)}$. A minimal
suffix there gives a minimal address in all of $S$, because immediate
children preserve the first letter. This proves the grafting assertion
using Lemma~\ref{lem:class_tree_grafting}.

First construct heights on individual address sets, without assuming
foundedness of the relation between trees. Founded recursion on the set
relation for $t$ gives the unique ordinal assignment
\[
r_t(p)=\sup\{r_t(q)+1:q\text{ is an immediate child of }p\}.
\]
This is the ordinal-rank construction following
Theorem~\ref{thm:setlike_class_recursion}; Replacement and ordinal union
give each supremum. Put $h(t)=r_t(\epsilon)$. Removing a subtree's
initial letter preserves its child relation. Transporting the restriction
of $r_t$ along $q\mapsto(a,b)q$ gives the subtree's rank assignment,
by uniqueness of founded recursion:
\[
r_{t_b}(q)=r_t((a,b)q).
\]
Hence
\[
h(t)=\sup_{b\in B(t(\epsilon))}(h(t_b)+1),\qquad h(t_b)<h(t).
\]
Each rank graph is a set by Replacement on its set domain $D_t$.
Consequently its existence and defining equations use only set
quantifiers. Elementary comprehension gives the class height
function $h:\mathcal W_B\to\On$.

The immediate subtrees of any $t$ form the image of the set
$B(t(\epsilon))$, so the subtree relation is set-like. If $S$ is any
nonempty subclass of $\mathcal W_B$, its ordinal height image is a
nonempty class of ordinals and has a least member. A tree in $S$ of
that height has no immediate subtree in $S$, since every such height
would be smaller. This proves foundedness on the class of trees.
\Done

\begin{theorem}[Initiality for class algebras with set arities]
\label{thm:class_W_initial}
For every class $X$ and class function $\alpha:F_B(X)\to X$, there
is exactly one class function $r:\mathcal W_B\to X$ satisfying
\[
r\circ\delta^{-1}=\alpha\circ F_B(r).
\]
An elementary predicate with fixed class parameters holds throughout
$\mathcal W_B$ if it is preserved by grafting from all immediate subtrees.
\end{theorem}
\Proof
Apply Theorem~\ref{thm:setlike_class_recursion} to the relation established
in Lemma~\ref{lem:class_W_founded}. At a tree $t=\delta^{-1}(a,s)$,
let $P_t$ be its set of immediate subtrees. For a set function $g$
with domain $P_t$ and range contained in $X$, Replacement on $B(a)$
forms $k(b)=g(s(b))$. Then $(a,k)\in F_B(X)$, so define the step
value to be $\alpha(a,k)$. On every other set input return the empty
set. This is a total set-valued class operation, defined by elementary
comprehension with the specified class parameters.

Founded recursion gives a class function $r:\mathcal W_B\to\V$.
Founded induction shows $r(t)\in X$: if this holds at all immediate
subtrees, the restricted predecessor graph is correctly typed, and the
step gives the value $\alpha(a,b\mapsto r(s(b)))\in X$. Thus the
default never occurs and the displayed morphism equation holds.
If two such class maps agree on the predecessors of a tree, their
equations give equal values there. A minimal disagreement is therefore
impossible, proving uniqueness. Repeated identical subtrees present no
ambiguity: the predecessor graph supplies one value for that subtree
at every occurrence.

For the induction assertion, elementary comprehension makes the class
of counterexamples. If nonempty, it has a minimal tree by the founded
subtree relation. All its immediate subtrees satisfy the predicate, so
grafting preservation contradicts its being a counterexample.
\Done

\subsection{Scope and a proper-class example}
\begin{example}[A proper class of nullary labels]
Take $A=\V$ and $B(a)=\emptyclass$ for every $a$. Every tree consists
of its root, with its chosen set label. Thus
\[
a\longmapsto\{(\epsilon,a)\}
\]
is a class bijection from $\V$ to both $\mathcal M_B$ and
$\mathcal W_B$. Each individual tree is a set. The carriers are proper
classes: if either were a set, Replacement by the root-label map would
make $\V$ a set. The initiality and finality statements above apply
to these proper carriers.
\end{example}

When $A$ is a set, Replacement collects the specified set fibres and
their tagged union is a set. The finite-address graphs and foundedness
condition then give exactly the preceding section's set constructions,
using the same pair and word coding. The class universal properties
also allow proper-class target algebras and source coalgebras; their
maps still have set values at each argument.

The extension therefore covers a proper class of possible set labels,
not individual labels that are themselves proper classes. It also
requires every arity to be a set. A proper-class branch fibre would
invalidate the set-level construction of the next address level and
the polynomial's set-function branch component. Theorems here make no
claim about that different interface, or about an encoding with
proper-class term values. Those versions still require their own
representation and recursion principles.

\ifdefined\qwertyflag
\undef\qwertyflag
\end{document}
\fi
